\documentclass[aps,pra,reprint,amsmath,amssymb]{revtex4-2}
\usepackage{graphicx}
\usepackage{xcolor}
\usepackage{bm}
\usepackage{tikz}
\usepackage{amssymb}
\usepackage{amsmath}
\usepackage{physics}
\usepackage{subfigure} 
% \usepackage[sorting=none]{biblatex}
% \addbibresource{References_local.bib}


\usetikzlibrary{shapes, positioning,decorations.pathmorphing, decorations.pathreplacing, arrows.meta}

\usepackage{quantikz}

\definecolor{bt}{RGB}{128, 0, 255}
\definecolor{erb}{RGB}{75, 0, 130}

\definecolor{mma1}{RGB}{83,118,172}
\definecolor{mma2}{RGB}{220,145,33}
\definecolor{mma3}{RGB}{132,167,44}
\definecolor{mma4}{RGB}{232,88,47}
\definecolor{mma5}{RGB}{124,109,170}
\definecolor{mma6}{RGB}{189,99,25}


% Redefine \section to produce inline bold text with optional TOC entry
\usepackage{titlesec}
\usepackage{hyperref}
% \titleformat{\section}[runin]  % inline/run-in heading
%   {\bfseries}                  % format for the heading text
%   {}                           % label (empty, so no numbering)
%   {0pt}                        % separation between label and title
%   {}                           % before-code (empty)

% \titlespacing*{\section}{0pt}{1.5ex plus .2ex minus .2ex}{1em}  % 

% % Redefine \subsection to be inline and italic
% \titleformat{\subsection}[runin]
%   {\itshape}
%   {}
%   {0pt}
%   {}

% \titlespacing*{\subsection}{0pt}{1ex plus .2ex minus .2ex}{1em}
\newcommand\perm[2][^n]{\prescript{#1\mkern-2.5mu}{}P_{#2}}
\newcommand\comb[2][^n]{\prescript{#1\mkern-0.5mu}{}C_{#2}}
\newcommand{\hsys}[0]{H_{\text{sys}}}
\newcommand{\henv}[0]{H_{\text{env}}}
\newcommand{\hint}[0]{H_{\text{int}}}
\newcommand{\UK}[0]{U_{\text{K}}}
\newcommand{\UI}[0]{U_{\text{I}}}
\newcommand{\Usys}[0]{U_{\text{sys}}[\bm{h}]}
\newcommand{\HK}[0]{H_{\text{K}}}
\newcommand{\HI}[0]{H_{\text{I}}}
\begin{document}

\title{
A Large N Correlated Noise Model on a Graph (Under Development)}

%Mode Protection and Synchronization of Anyonic Oscillators}
\author{BT}
\email{btyagi@uh.edu}
\affiliation{Department of Physics, University of Houston, Houston, Texas, 77204, USA}
\author{ERB}
\affiliation{Department of Physics, University of Houston, Houston, Texas, 77204, USA}
\author{KEB}
\affiliation{Department of Physics, University of Houston, Houston, Texas, 77204, USA}
\date{\today}


\begin{abstract}
%We derive the Lindblad master equation for anyonic oscillators and extend it to coupled modes under correlated noise. Fractional statistics modify normal-mode structure and interpolate between bosonic and fermionic dissipation. By tuning the anyon phase and environmental correlations, we achieve selective mode protection and observe exceptional points in the dissipative spectrum. These results show that exchange statistics enable tunable coherence control in topological quantum systems.
We generalize the correlated noise model to a large system by identifying the correlation matrix as the topological adjacency of a connected graph. We want to see how the dynamics change as the system becomes increasingly dense as we increase the edges or the connectivity of the graph. We want the see the effects of the topology of the graph and how does the degree of correlation affect the spread of information in the system by analysing the behavior of OTOCs. Finally we also want to strongly characterize the amount of entanglement in the system by calculating the entanglement negativity. 
\end{abstract}



\maketitle

\section{Introduction}
A key insight into the dynamical properties of physical systems is given by how bits of information scramble as the system evolves in time. According to Landauer's Principle \cite{landauer1961irreversibility}, this information is related to the systems entropy and is indeed physical, even in quantum systems \cite{peterson2016experimental}, and therefore understanding how the mechanics is manifest in information transport is a very challenging, yet interesting problem in local quantum systems. A myriad of pivotal, numerical and analytical, results exist for isolated quantum systems that evolve under unitary dynamics. However, most physical systems of nature, on top of being {\it non-integrable} and {\it chaotic} are also constantly interacting with an environment whose degrees of freedom we can not fully account for. Such systems are famously called {\it open} quantum systems \cite{breuer2002theory}. It is this interplay between internal chaotic dynamics and the constant interaction with the environment that makes the study of these systems challenging but highly rewarding since such systems form the backbone of most experimental setups \cite{bergamasco2023quantum}. While there have been studies that try to weakly couple closed systems to some sort of environment and study the effects of dissipation and/or pumping on the dynamics of these systems, all such studies treat noise as an enemy of \textit{quantum coherence} and attempt to shield the system from its effects thereby separating the effects of scrambling and decoherence \cite{zhang2019information}. However, the role of the environment, its geometrical structure and correlations within this environment, while noting its importance, is largely neglected. In our recent work \cite{Tyagi:2024aa,Bittner2024,bittner2025statistical}, we flip this narrative by showing that {\it quantum coherence} can actually be preserved for long timescales by treating this noise as an ally, if we factor in the above missing crucial ingredients. A widely studied measure of information scrambling, also related to the quantum Lyapunov exponent, is the out-of-time ordered correlator (OTOC)\cite{xu2024scrambling}. The OTOC is intuitive for a closed system while what would be the properties of this signal in open systems is still an open problem. In physical systems, local environments or noise-channels tend to be {\it spatially correlated}, such noise-correlations can lead to interesting dynamical features.  Our system, weakly coupled to an environment,  evolves via the {\it Gorini, Kossakowski, Sudarshan and Lindblad} (GKSL) master equation given by 
\begin{align}
\label{GKSL}
    \frac{d\rho}{dt} =-i [H_{\text{sys}},\rho]+ \sum_k \gamma\left[ L_k \rho L_k^\dagger - \frac{1}{2}\{\rho,L_k^\dagger L_k\}\right].
\end{align}
This can be compactly written as 
$\frac{d\rho}{dt} = \mathcal{L}(\rho)$ where $\mathcal{L}$ is the {\it Liouvillian} superoperator which generates the evolution corresponding to a completely-positive-trace-preserving (CPTP) \cite{lindblad1976generators} dynamical map $\mathcal{M}(t):\rho(0) \rightarrow \rho(t)$.
Where $\gamma$ is the overall dissipation and operators $L_k$ are the local Lindblad operators.
%are defined using the eigenvalues and eigenvectors of the correlation matrix
% \begin{align}
%     L^{(em)}_\alpha&= \sqrt{\gamma (\bar{n}+1)\tilde{\lambda}_\alpha}\tilde{a}_\alpha\\
%     L^{(abs)}_\alpha &= \sqrt{\gamma \bar{n}\tilde{\lambda}_\alpha}\tilde{a}^\dagger_{\alpha}
% \end{align}
% where $(\tilde{a}_\alpha,\tilde{a}^\dagger_\alpha)$ are transformed bosonic operators $\tilde{a}_\alpha(\tilde{a}_\alpha^\dagger)=\sum_{i=0}^{N-1}T_{i \alpha}a_\alpha(a^\dagger_\alpha)$. 
In the Heisenberg picture 
%where the operators evolve with time,
where a general local operator $O$ has the expectation value
$\langle O(t) \rangle = tr\left(\rho O(t) \right)$, taking the time-derivative we have 
\begin{align}
    \frac{d \langle O_H \rangle}{dt}
%&=tr\left(\dot{\rho}O_H\right)=tr\left(\rho \dot{O_H}\right)\nonumber\\
    % &=tr\left(\mathcal{L}(\rho)O_H\right)=tr\left(\rho \mathcal{L}^\dagger (O_H)\right) \nonumber \\
    % &
    = &\langle i[H_{\text{sys}},O_H] \rangle + \nonumber\\
    &\langle \sum_k \gamma \left[ L^\dagger O_H L - \frac{1}2{}\{O_H,L^\dagger_k L_k\}\right]\rangle
\end{align}
where $\mathcal{L}^\dagger$ is now the {\it adjoint Liouvillian} superoperator generating the the adjoint dynamical map $\mathcal{M}^\dagger(t):O_H(0)\rightarrow O_H(t)$. For simple correlated noise-channels the non-unitary part in Eq. \eqref{GKSL} splits into a symmetric and an anti-symmetric part and as a consequence only affects certain modes of the system while the others are completely immune to this environmental noise and our construction allows for complete mode selectivity. We introduce a control parameter $\xi$ and show that at certain critical values, our system admits a {\it dynamical phase transition}, it {\it spontaneously synchronizes}, similar to a Huygens' clock \cite{huygens1999instructions} and remains perpetually locked in a non-equilibrium steady state (NESS).  Moreover we see that even in the absence of a direct coupling between two-quantum systems, the noisy environment can induce resourceful quantum correlations like {\it entanglement} in the system. 

In this letter we show that correlated noise can dramatically alter information scrambling in the quantum systems.  

\section{Correlated Noise Model}
% \section{System 1: 3D Anderson Model}
% The tight-binding 3D Anderson Model on a cubic lattice given by
% % \begin{align}
% % \label{Hamiltonian}
% %     H_{\text{sys}} &= H_0 + H^{(m)}_{\text{int}} \nonumber\\
% %     H_{\text{sys}} &= \sum_{i=0}^{N-1} \omega_i a_i^\dagger a_i +\sum_{(i,j)\in \mathcal{E}_m} J \left(a^\dagger_i a_j + a^\dagger_j a_i\right)
% % \end{align}
% % where $H_0$ represents uncoupled harmonic oscillators and $H_{\text{int}}^{(m)}$ represents the interaction Hamiltonian for a graph with $m$ edges where $m=N,N+1, N+2,\cdots, \comb[N]{2}$.

% % $\{\omega_i\}$ is the frequency of each local mode, $(a^\dagger_i,a_i)$ are creation and annihilation operators for mode $i$, obey the usual canonical commutation relations $[a_i,a^\dagger_j]=\delta_{ij}$, $J$ is the interaction strength and $\hbar=1$. The sets are defined as
% % \begin{align*}
% %     \mathcal{E}_m &= \{(i,j)\in \{0,1,\cdots (N-1)\}^2 | i < j, A_{ij}^{(m)}=1\}\\
% %     \mathcal{R}_k &= \{(i,j)\in \{0,1,\cdots (N-1)\}^2 | i < j, A_{ij}^{(k)}=1\} 
% % \end{align*}
% % where $\mathcal{E}_m$ is the general set of all graphs with $m$ edges with adjacency matrix given by $\mathbf{A}^{(m)}$ and $\mathcal{R}_k$ is the set of $k-$regular graphs for $N$ nodes with the regular adjacency matrix $\mathbf{A}^{(k)}$ and $k=2,3,\cdots,N-1.$ and we know 
% % \begin{align*}
% %     \mathcal{E}_m =\begin{cases}
% %         \mathcal{R}_k \text{ if }m=kN/2. \\
% %         non-regular \text{ otherwise.}
% %     \end{cases}
% % \end{align*}
% % The Hamiltonian (\ref{Hamiltonian}) can be compactly written as 
% \begin{align}
%     \hsys &=  - J \sum_{i\neq j} A_{ij}(c^\dagger_i c_j+h.c) +\frac{W}{2}\sum_{i=1}^{N} \omega_i c_i^\dagger c_i
% \end{align}
% where $N$ is the number of spins, $\{\omega_i\}$ are independent and identically distributed (i.i.d) random numbers we draw from a uniform distribution from the interval $[-1,+1]$ and $W$ corresponds to the disorder amplitude. This model is equivalent to a simple tight-binding model with random on-site energies and is a type of {\it single-particle} quantum-chaotic quadratic Hamiltonian system \cite{lydzba2021single, lydzba2021entanglement} below localization transition, which in this case is around $W_c \approx 16.5$ \cite{slevin2018critical}. Restricting to at most one excitation in the system, the total Hilbert space is the direct sum $\mathcal{H}=\mathcal{H}_{\text{vac}}\oplus_{i=1}^N\mathcal{H}_{\text{sys}_i}$. This model has been widely studied for its transport and spectral properties \cite{vsuntajs2021spectral}. 
We consider the \textit{dissipative} dynamics of a minimal chaotic model \cite{prosen2002general,prosen2002ruelle,prosen2007chaos} with local interactions and a finite dimensional local Hilbert space. The chaotic properties of this model are described by comparing spectral properties with those of Random Matrices \cite{pineda2007universal}. We turn on correlated dissipation using a graph theoretic structure and account for how entanglement dynamics and scrambling rates change in the weak dissipation limit. We show that the graph structure protects some modes of the system from environmental noise. 
\section*{System: Kicked Field Ising Model (KFIM)} 
\begin{align}
\label{SysHam}
\hsys[{\mathbf h};t] &= H_{\text{I}}[\bm{h}]+\sum_{m=-\infty}^{\infty} \delta(t-m) H_{\text{K}}\\
H_{\text{I}}&\equiv \frac{J}{2}\sum_{i, j=1}^L A_{ij}\sigma^z_i \sigma^z_{j}+\sum_{j=1}^L h_j \sigma^z_j, \nonumber\\
H_{\text{K}}&\equiv b\sum_{j=1}^L \sigma^x_j. \nonumber
\end{align}
In this ``self-dual" 1D Floquet-Ising spin chain with periodic boundary conditions, $L$ represents the system volume and system operators $\sigma^a_j$ are the Pauli matrices, with $a\in\{x,y,z\}$, defined at position $j$. System parameters $J,b$ and ${\mathbf h}=(h_1,\cdots,h_L)$ correspond to the Ising coupling, the transverse kicking field and the position dependent longitudinal field, respectively.
$\mathbf{A}$ is the {\it adjacency} matrix of a given graph such that 
\begin{align*}
A_{ij}=
    \begin{cases}
       1 \text{ if $i \neq j$ $\leftrightarrow$  there's an edge}\\
       0 \text{ otherwise }.
    \end{cases}
\end{align*}
The non-integrability of this model, by calculating exactly, the level-statistics, was proven by Bertini, Kos and Prosen \cite{bertini2018exact} who also showed the existence of a special class of initial states, from which entanglement dynamics can be computed exactly, making this a minimal locally interacting model, with a finite dimensional local Hilbert space that exhibits quantum chaos \cite{bertini2019entanglement}. Note that for ${\mathbf h}=0$ this model is integrable and is chaotic for any non-zero value of the longitudnal fields. In the parameter space of this model, we work at these ``self-dual" points that give us a strong analytic handle on the dynamics and so we set $J=\pi/4$ and $b=-\pi/4$ from here on. We start with the initial state 
\begin{align}
    \rho(0)&=\ket{\psi_{\bm{\theta},\bm{\phi}}}\bra{\psi_{\bm{\theta},\bm{\phi}}} \nonumber \\
    \ket{\psi_{\bm{\theta},\bm{\phi}}}&= \bigotimes_{j=1}^L \left[\cos \left(\frac{\theta_j}{2}\right)\ket{\uparrow}+\sin \left(\frac{\theta_j}{2}\right)e^{i \phi_j}\ket{\downarrow}\right].
\end{align}
The special class of states also known as separating states are defined as
\begin{align}
    \mathcal{T}&=\left\{\ket{\psi_{\frac{\pi}{2}\bm{1}.\bm{\phi}}}|\;\phi_j \in [0,2\pi]\right\} \\
    \mathcal{L}&= \left\{\ket{\psi_{\bar{\bm{\theta}},\bm{\phi}}}|\;\bar{\theta}_j \in \{0,2\pi\}\right\}
\end{align}
where $\mathcal{T}$ and $\mathcal{L}$ are called transverse and longitudinal separating states respectively. The time-evolution generated by (\ref{SysHam}) is given by the following one-period Floquet operator
\begin{align}
\label{Flqop}
    U_{\text{sys}}[\mathbf{h}]&= T\exp \left[-i\int_0^1 dt \hsys[\bm{h};t]\right]\nonumber\\
    &=e^{-i\HK}e^{-i\HI[\bm{h}]}\\
    &\equiv \UK\UI \nonumber.
\end{align}
The reduced density matrix for an arbitrary contiguous bipartition of the system consisting of $N$ spins is given by 
\begin{align}
    \rho_A(t)=\Tr_{\mathcal{H}_{L-N}} \rho(0)
\end{align}
where $A=\{1,2\cdots,N\}$ and $\mathcal{H}_{L-N}$ is the Hilbert space of the compliment $A^c=\{N+1,\cdots,L\}$. The entanglement entropy of the subspace described by $\rho_A(t)$ is given by the Rényi entropies 
\begin{align}
    S_A^{(\alpha)}(t)=\frac{1}{1-\alpha}\log \Tr \left[(\rho_A(t))^{\alpha}\right],\quad \alpha >0.
\end{align}
The time-evolution, in the thermodynamic limit $(L \rightarrow \infty)$, from states in the separating class $\mathcal{T}$ has been explicitly shown to take the following universal form
\begin{align}
\label{universalEE}
    \lim_{L \rightarrow \infty} S_A^{\alpha}(t)= \min(2t,N)\log 2.
\end{align}
See Appendix \ref{App:A} for all relevant details of the model.
\section*{Environment: Noise-Correlation Matrix}
As discussed in \cite{Bittner:2025aa}, the noise-correlation matrix ${\bf \Xi}$, that encodes the correlations between noise channels (environment) can describe the topological adjacency between the correlated Wiener processes $dW_i(t)dW_j(t') = \Xi_{ij}\delta_{ij}\delta(t-t')$ and can be constructed from a connected- cyclic graph belonging to a particular permutation group. The eigenvectors of the correlation matrix are then given by the irreducible representations of this group. Note that ${\bf \Xi}$ is required to be positive semi-definite to have a well defined quantum channel. For a $k$-regular graph, the correlation matrix ${\bf \Xi}$ can take the following minimal ansatz
\begin{align}
    \mathbf{\Xi}=\mathbf{I}-\xi \mathbf{A}^{(k)} 
\end{align}
where $\mathbf{A}^{(k)}$ is the {\it Adjacency} matrix of the graph. 
% For the case of two-oscillators this was 
% \begin{align}
%     \mathbf{\Xi}=\begin{pmatrix}
%         1 & \xi \\
%         \xi & 1
%     \end{pmatrix}
% \end{align}
% For a general graph it is suitable to work with the {\it Laplacian} Matrix which typically has the degree of the graph on the diagonal and 
% \begin{align*}
% L_{ij}=
%     \begin{cases}
%         -1 \text{ if there's an edge}\\
%         0 \text{ otherwise}
%     \end{cases}
% \end{align*}
% For $k$-regular graphs the {\it Laplacian} and the {\it Adjacency} matrix are closely related 
% \begin{align}
%     \mathbf{L} = k \mathbf{I}-\mathbf{A}^{(k)}
% \end{align}
% If $v$ is an eigenvector of $\mathbf{A}^{(k)}$ with eigenvalue $\lambda$ then one can show that $v$ is also an eigenvector of $\mathbf{L}$ with eigenvalue $k-\lambda$. The {\it Laplacian} matrix admits a quadratic form 
% \begin{align}
%     \mathbf{x}^T \mathbf{L} \mathbf{x} = \sum_{i,j} (x_i - x_j)^2
% \end{align}
% since this is a sum of squares it is a proof that the {\it Laplacian} matrix for any graph is positive semi-definite. 
% Therefore we can transform our correlation matrix into the {\it Laplacian} form
% \begin{align}
%     \mathbf{\tilde{\Xi}}&= k \mathbf{\Xi} \nonumber \\
%     &=k \mathbf{I} - \tilde{\xi} \mathbf{A}^{(k)}
% \end{align}
% where $\tilde{\xi} = k \xi$ is the modified correlation parameter. In order to maintain the positive semi-definite property it is easy to show that 
% \begin{align}
%     \xi &\in \left[\frac{-1}{k},\frac{1}{k}\right]
% \end{align}
% where $\xi = -1/k$ corresponds to \textit{anti-correlated noise}, $\xi=1/k$ corresponds to \textit{correlated noise} and $\xi=0$ corresponds to \textit{uncorrelated noise}.
% Note that in our setup the noise-correlation matrix will always be {\it regular} and can be independently constructed which means that it doesn't have to strictly depend on the system Hamiltonian. As we will show later that we can make the system graph more and more dense without changing the environment structure.
Since $\mathbf{\Xi}$ is positive semi-definite, real, symmetric matrix we can diagonalize it 
% $\mathbf{T}^\dagger \mathbf{\Xi}\mathbf{T} = diag\{\lambda_\alpha\}$
and we know that for any cyclic graph $C_\alpha$, the eigenvalues of the adjacency matrix are given by $c_\alpha = 2\xi\cos{k_\alpha}$ such that the discrete Fourier modes are $v_\alpha(j)= \frac{1}{\sqrt{L}}e^{ik_\alpha j}$ where $k_\alpha=\frac{2\pi \alpha}{L}$, $L$ is the number of vertices or Lattice sites and $\alpha=0,1,\cdots,n-1$. Then for the correlation matrix $\bf{\Xi}$, the eigenvalues are $\lambda^{\Xi}_\alpha(\xi)=1+\xi \lambda_\alpha$. Since $\cos{k_\alpha}\in [-1,1]$, $\lambda^{\Xi}_{\text{min}}=1 - 2|\xi|$. For the noise-correlation matrix to be positive semi-definite $\xi \in [-1/2,1/2]$. 
In the correlated dissipation analysis we will focus on two site-wise dissipative quantum channels
\\[4pt]
1. Dephasing:
\begin{align}
    L_\alpha =  \sqrt{\gamma \lambda_\alpha} \sum_{j=1}^Lv_{\alpha}(j)\sigma^z_j
\end{align}
2. Amplitude:
\begin{align}
    L_\alpha^- &= \sqrt{\gamma (n_{\text{th}}+1)\lambda_\alpha}\sum_{j=1}^L v_\alpha(j)\sigma^-_j \nonumber\\
    L^+_\alpha &= \sqrt{\gamma (n_{\text{th}})\lambda_\alpha}\sum_{j=1}^L v_\alpha(j)\sigma_j^+
\end{align}
where $\{\lambda_\alpha,v_\alpha\}$ are eigenpairs of the correlation matrix $\Xi$, $\sigma_j^{\pm}=(\sigma^x_j \pm \sigma^y_j)/2$, and $n_{\text{th}}(\omega)$ is the Bose-Einstein distribution
\begin{align}
    n_{\text{th}}(\omega) \equiv \frac{1}{e^{\beta\omega\hbar} -1}
\end{align}
where $\omega$ is the transition frequency $\omega$. Note that since pure dephasing only acts in the coherence subspace and doesn't change populations in the system the effect of bath induced entanglement will be manifest in the amplitude damping channel where the system generally evolves towards a non-trivial steady state. Physically interesting are special points for example at $\xi=-1/2$ (anti-correlated noise), $$\lambda_\alpha=1-\cos{k_\alpha},\quad \lambda_0=0.$$ This corresponds to the dark (or dissipation free) mode
\begin{align*}
    v_0 (j)=\frac{1}{\sqrt{L}}.
\end{align*}
So the corresponding collective dephasing channel \begin{align*}
    L_0 = \sqrt{\gamma \lambda_0} \sum_{j=1}^L\frac{1}{\sqrt{L}}\sigma^z_j = 0.
\end{align*}
and for $\xi=1/2$ (correlated noise) the dark mode corresponds to $k_\alpha=\pi$ for $\alpha=L/2$ (assuming even L)
\begin{align*}
    v_{L/2}(j)= \frac{1}{\sqrt{L}}(-1)^j
\end{align*}
and so the corresponding dephasing channel
\begin{align*}
    L_{L/2} = \sqrt{\gamma \lambda_0} \sum_{j=1}^L\frac{(-1)^j}{\sqrt{L}}\sigma^z_j = 0.
\end{align*}
Similarly for the Amplitude damping channel
% \subsubsection{Spectral Analysis}
\section{Dynamics}
For the Open Floquet System, in the weak dissipation limit $(\gamma \ll 1)$, we can define two equivalent evolution scenarios for dissipation up to some $\mathcal{O}(\gamma)$ corrections.
\subsection{Continuous Channel}
Notice that the instantaneous kicks in the evolution generated by (\ref{Flqop}) take place at integer steps and we can allow for continuous dephasing between the kicks. For $t \in (m,m+1)$ we have 
\begin{align}
    \dot{\rho} = -i[H_\text{I}, \rho] + \mathcal{D}(\rho) \equiv \mathcal{L}_{\text{cont}}(\rho).
\end{align}
% Then at $t=m$, we apply the kick
% \begin{align}
% \rho(m^+)=\UK\rho(m^-)\UK^\dagger, \quad \UK=e^{-i \HK}
% \end{align}
% Therefore the dynamical map between an integer period is given by 
% \begin{align}
%     \rho_{m+1}=\UK \exp^{\mathcal{L}_\text{cont}}(\rho_m)\UK^\dagger .
% \end{align}
where $\mathcal{L}_{\text{cont}}$ is the continuous channel Liouvillian. We can therefore define the continuous channel dynamical map as the composition of continuous time dissipation followed by a kick 
\begin{align}
    \mathcal{M}_{\text{cont}}(\rho):=\mathcal{U}_\text{K} \circ e^{\mathcal{L}_{\text{cont}}}(\rho).
\end{align}
\subsection{Stroboscopic Channel}
For the stroboscopic channel we simply allow for the dissipation to accumulate over the integer time-period and do not resolve the micro-time dissipative dynamics within that period and so the system density matrix evolves via the following map
\begin{align}
    \mathcal{U}_{\text{sys}}(\rho)&= U_{\text{sys}}[\bm{h}] \rho U^\dagger_{\text{sys}}[\bm{h}] \nonumber\\
\end{align}
We then introduce a dissipation channel $\mathcal{E}_{\text{diss}}$ that acts once per period, and the complete dynamical map looks like
\begin{align}
\label{strobdynamics}
    \mathcal{M}_{\text{strob}}(\rho):=\mathcal{E
    }_{\text{diss}} \circ \mathcal{U}_{\text{sys}}(\rho).
\end{align}
Such a stroboscopic channel closely resembles experimental setups, where drives are practically instantaneous and periodic and one only cares about integer stroboscopic times where Markovianity ensures that the environment is featureless on the scale of a period. So this can be thought of as a coarse-grained Floquet quantum channel.
\\[4pt]
In the weak dissipation limit $(\gamma \ll 1)$ both the continuous and the stroboscopic channel are equivalent up to $\mathcal{O}(\gamma)$ corrections. See Appendix \ref{App:C}. 
\subsection{Out of Time Ordered Correlators (OTOCs)}
%  Larkin and Ovchinnikov \cite{larkin1969quasiclassical}  introduced Out-of-Time-Ordered Correlators (OTOCs) in the context of superconductors. They showed that
% \begin{align}
%     e^{\lambda t} \sim\{q(t),p\}_{PB} = \lim_{\hbar \rightarrow 0} [\hat{q}(t),\hat{p}]
% \end{align}
% which suggests that {\it Quantum Lyapunov} hides in $[\hat{q}(t),p]$.
The Out-of-time-ordered Correlator (OTOC) is a robust measure of how information scrambles in a system and is manifest in the growth of the norm of the commutator $\left[W(t),V(0)\right]$ where $W$ and $V$ are local operators\cite{larkin1969quasiclassical,xu2024scrambling}.
% \begin{align}
%     \langle [q(t),p]^2\rangle= \nonumber 
%     \langle{q_t p_0 q_t p_0 + p_0 q_t p_0 q_t+\text{time-ordered terms}}\rangle.
% \end{align}
For Hermitian and Unitary operators, generally, the true OTOC takes the following nice form. See Appendix \ref{App:A}
\begin{align}
    F(t,r)&=\frac{1}{d}\Re\langle  W^\dagger(t)V^\dagger(0)W(t)V(0)\rangle_\beta \nonumber\\
    &=1-\frac{1}{2d}\langle[W(t),V(0)]^\dagger[W(t),V(0)]\rangle_\beta
\end{align}
where $\langle \cdot \rangle_{\beta}$ represents the thermal average at temperature $\beta = 1/k_BT$ and $d$ is the Hilbert space dimension. 
% \begin{align}
%     C(t,r)=\langle{\underbrace{[W(t,0),V(0,r)]^\dagger [W(t,0),V(0,r)]}_{\text{norm of the commutator}}}\rangle.
% \end{align}
% The full OTOC is then
% \begin{align}
%     C(t,r) = -F(t,r)-F^*(t,r)+\underbrace{\text{time-ordered}}_{\text{decays by }t_{\beta}\text{ ``dissipation time"}}.
% \end{align}
Figure (\ref{fig:OTOC}) shows that at some initial time $t_0$ we have two system operators $W(0,0)$ and $V(0,r)$.
\begin{widetext}
\begin{figure}[htpb]
    \centering
    \includegraphics[width=\textwidth]{Figures/otoc_spread.png}
    \caption{Caption}
    \label{fig:OTOC}
\end{figure}
\end{widetext}
Initially, their supports have no overlap and therefore their commutator vanishes. What that means is that the operators are essentially independent and act on separate parts of the system's Hilbert space. They do not interfere with each other, so their order of application does not matter. With time the operator $W(t)$ grows under Heinsenberg evolution 
\begin{align}
    W(t,0) = e^{iHt}We^{-iHt}
\end{align}
and its support increases. The squared commutator grows as the influence of $W(t)$ reaches $V$.
% \begin{figure}[h]
%     \centering
%     \includegraphics[width=0.9\linewidth]{Pictures/OTOC.png}
%     \caption{Pictorial representation of OTOC and operator growth}
%     \label{fig:OTOC}
% \end{figure}
% So now what exactly do we mean by "operator size" and "complexity"?
In a $q-$local chaotic system, we have 
\begin{align}
    C(t,r) \sim \frac{1}{N} e^{2 \lambda(t-r/v_B)}
\end{align}
where $\lambda$ corresponds to the {\it Quantum Lyapunov Exponent}. This behavior of the OTOC is seen for large$-N$ or in the semi-classical limit where the scrambling time $t_* \sim \ln N$. Maldacena-Shenker-Stanford (MSS) show a bound for the quantum lyapunov exponent or generally a bound on chaos \cite{maldacena2016bound}
\begin{align*}
    \lambda \leq 2 \pi T.
\end{align*}
% \subsubsection{Dissipative Scrambling}
For open systems, calculating the OTOC is complicated because the system constantly interacts with its environment and suffers dissipation. Since the system evolves via the Lindblad Master equation, where time-reversal symmetry is broken by construction and the Born-Markov-Secular approximation is assumed, there is no obvious way to calculate OTOCs in real-world experiments since they are inherently time-ordered and also that any information scrambling signal would inevitably decay due to dissipation. A possible way out and the method we adopt in this paper is to study the structural changes in the scrambling light cone to characterize scrambling \cite{zhang2019information}. A dissipative system is coupled to an environment whose total Hamiltonian is given by
\begin{align*}
    H=\hsys+\henv+\hint
\end{align*}
corresponding to the Hamiltonian of the system, environment and the interaction term respectively. So in the presence of dissipation, we study the OTOC signal following a robust measurement scheme in \cite{swingle2016measuring}, applied in \cite{zhang2019information}, for a Chaotic system, and further improved in \cite{swingle2018resilience}. Since the forward-backward time ordering of operators in the OTOC forces us to think about the dynamics carefully, especially in the presence of dissipation. Naively one would think that $-H=-(\hsys+\henv+\hint)$ would give us time reversed evolution but that is simply not true because in an experiment we have no control over $\henv$ and $\hint$. So practically, only $\hsys \rightarrow -\hsys$ corresponds to backward time-evolution. So we can define $H_f=\hsys + \henv + \hint$ which controls forward evolution and $H_b = -\hsys +\henv +\hint$ which controls the backward time evolution. Therefore the dynamical and the adjoint dynamical maps defined above also split into a forward and backward component, where the backward maps come with a minus sign in front of the system Hamiltonian.
\section{Results}
We pick a state in the transverse class-$\mathcal{T}$ where each spin is an eigenstate of an operator in the xy-plane, then for $\phi=0$ and $\theta=\pi/2$ we have 
\begin{align}
    \ket{\psi}=\bigotimes_{j=1}^L \ket{+}_j,\quad\ket{+}=\frac{1}{\sqrt{2}}\left(\ket{\uparrow}+\ket{\downarrow}\right).
\end{align}
Although the dephasing acts continuously during each Floquet period, we sample the dynamics stroboscopically at integer times and describe it via the Stroboscopic channel. So we write down one Floquet step with dissipations as in (\ref{strobdynamics}) and
for weak dissipation we have $\gamma\ll1$, 
\begin{align}
    \rho_{t}&= \Usys \rho_0\Usys^\dagger \nonumber\\
    &\quad+ \mathcal{D_{\text{diss}}}\left(\Usys \rho_0\Usys^\dagger\right).
\end{align}
For notational simplicity let $$X_t:=\rho_t^{T_{A^c}}$$ be the partially transposed state, since the partial transpose operation is linear
\begin{align}
\label{parttranspo}
    X_{t+1}&= \left(\Usys \rho_t\Usys^\dagger\right)^{T_{A^c}} \nonumber\\&\quad+ \left(\mathcal{D_{\text{diss}}}\left(\Usys \rho_t\Usys^\dagger\right)\right)^{T_{A^c}}.
\end{align}
Now lets consider this evolution for different channels
\subsubsection{Dephasing}
Since Pauli operators are hermitian $(\sigma)^\dagger =\sigma$ we can write the dephasing channel as 
\begin{align}
    \mathcal{D}_{\text{deph}}(\rho)=\sum_\alpha \gamma \lambda_\alpha \left(Z_\alpha\rho Z_\alpha-\frac{1}{2}\{\rho,Z_\alpha^2\}\right)
\end{align}
where $Z_\alpha := \sum_{j=1}^L v_\alpha(j)\sigma^z_j$.
Additionally dephasing operator is the sum of diagonal matrices $\sigma_z$, they're unchanged by the partial transpose, so
\begin{align}
\mathcal{D}_{\text{deph}}^{T_{A^c}}(\rho)= \sum_\alpha \gamma \lambda_\alpha \left(Z_\alpha\rho^{T_{A^c}} Z_\alpha-\frac{1}{2}\{\rho^{T_{A^c}},Z_\alpha^2\}\right) 
\end{align}
and therefore the dephasing evolution of the partially transposed state becomes
\begin{widetext}
\begin{align}
    X_{t+1}= \left(\Usys\rho_t \Usys^{\dagger}\right)^{T_{A^c}}
    +\sum_\alpha \gamma \lambda_\alpha \left(Z_\alpha\left(\Usys\rho_t\Usys^\dagger\right)^{T_{A^c}} Z_\alpha-\frac{1}{2}\{\left(\Usys\rho_t\Usys^\dagger\right)^{T_{A^c}},Z_\alpha^2\}\right).
\end{align}
\end{widetext}
Under weak dissipation 
\begin{align}
    X_{t+1}&= \left(\mathbb{I}+\gamma\mathcal{D}_{\text{deph}}\right)\left(\Usys\rho_t\Usys^\dagger\right)^{T_{A^c}}+O(\gamma^2).
\end{align}
\subsection{Logarithmic Negativity}
For {\it open} systems Rényi entropies are not a robust entanglement measure since the global state of the system is mixed. In other words, the ``mixedness" of some subsystem is not only because of quantum correlations but also due to the interaction with the environment. So instead one can calculate the {\it logarithmic negativity} (LN) which is defined by
\begin{align}
    \mathcal{E}(\rho;A):=\ln \Tr{|\rho^{T_{A^c}}|}, 
\end{align}
where $T_{A^c}$ is the partial transpose on subsystem $A^c:=\{N+1,\cdots,L\}$. When dealing with pure, separating states as in \cite{bertini2019entanglement}, LN simply is equal to the Rényi entropy where $\alpha=1/2$ (See Appendix \ref{App:D}). Therefore the unitary LN for states belonging to class-$\mathcal{T}$ obey the same universal law in the thermodynamic limit (\ref{universalEE})
\begin{align}
    \mathcal{E}_{\text{unitary}}(t)=\min(2t,N)\ln2.
\end{align}
We can treat this as our analytic baseline to compare what happens when we turn on localized dissipation. 
% where $\Usys^{T_{A^c}}$ is the partial transpose of the unitary matrix $\Usys$.
% \begin{figure*}
%     \centering
%     \includegraphics[width=0.5\linewidth]{Figures/noisy_graph_model.pdf}
%     \caption{Caption}
%     \label{fig:adding edge}
% \end{figure*}
\subsubsection{Amplitude Damping}
\begin{align}
    X_{t+1}&= \left(\Usys\rho_t \Usys^{\dagger}\right)^{T_{A^c}}\nonumber\\
    &\quad+\left(\mathcal{D_{\text{amp}}}\left(\Usys \rho_t\Usys^\dagger\right)\right)^{T_{A^c}}.
\end{align}
\begin{acknowledgments}
This work was supported by the National Science Foundation (CHE-2404788) and the Robert A. Welch Foundation (E-1337). 
\end{acknowledgments}



\section*{Data Availability Statement}

\section*{Competing Interests}

The authors declare no competing interests.

\section*{Author Contributions}


\bibliographystyle{apsrev4-2}  % Choose your preferred style
\bibliography{References_local}    % Your .bib file (without the .bib extension)

\newpage
\appendix

\begin{widetext}

\section{Kicked Field Ising Model}
\label{App:A}
% \subsection{3D Anderson Model}
% The Hilbert space can be written as the direct sum
% \begin{align}
%     \mathcal{H}=\mathcal{H}_{\text{vac}} \oplus \bigoplus_i^N \mathcal{H}_{\text{sys}_i}
% \end{align}
% where $N$ is the number of degrees of freedom, $\mathcal{H}_{\text{vac}}=span\{\ket{0}\}$ and $\mathcal{H}_i = span\{\ket{i}\}$. We work in the orthonormal basis $\ket{0},\ket{1},\cdots \ket{N}$ where $\ket{0}$ corresponds to the vacuum state and $\ket{i}$ represents the state where the excitation is localized at site $i$. The Hamiltonian of the system is
% \begin{align}
%     H_{\text{sys}} &=  - J \sum_{i\neq j} A_{ij}(c^\dagger_i c_j+h.c.) +\frac{W}{2}\sum_{i=1}^{N} \omega_i c_i^\dagger c_i
% \end{align}
% where, in the restricted subspace, the creation, annihilation and the density operators $c^\dagger_i =\ket{i}\bra{0}$, $c_i=\ket{0}\bra{i}$ and $n_i=c^\dagger_i c_i=\ket{i}\bra{i}$, obey the following relations
% \begin{align*}
% a^\dagger_ia^\dagger_j=a_ia_j=0,[n_i,a_j^\dagger]=\delta_{ij}a_i^\dagger,\;[n_i,a_j]=-\delta_{ij}a_i
% \end{align*}
% Note that in the single-particle sector, the statistics don't matter. So we have 
% \begin{align}
% H\ket{0}=0,\;    
% \end{align}
% \subsection{Mixed Field Ising Model}
\section{OTOCs}
\label{App:B}
The squared norm of the commutator is given by
\begin{align}
    C(t)&=[W(t),V]^\dagger [W(t),V]\nonumber\\
    &=\left(V^\dagger W^\dagger(t)-W^\dagger(t) V^\dagger\right)\left(W(t)V-VW(t)\right)\nonumber\\
    &=V^\dagger W^\dagger(t) W(t)V-V^\dagger W^\dagger(t)V W(t)-\nonumber\\&W^\dagger(t)V^\dagger W(t)V+W^\dagger(t)V^\dagger V W(t)
\end{align}
Let $W$ and $V$ be hermitian and unitary. So $W=W^\dagger$, $V=V^\dagger$ and $W^\dagger W=V^\dagger V=\mathbf{1}$. Then taking the trace we get
\begin{align}
    Tr(C(t))&=2(Tr(\mathbf{1})-\Re Tr(W^\dagger(t) V^\dagger W(t) V))\nonumber\\
    \frac{1}{2d}Tr(C(t))&=1-\frac{1}{d} \Re Tr(W^\dagger(t) V^\dagger W(t) V)
\end{align}
The OTOC is $F(t)=\frac{1}{d}\Re Tr(W^\dagger(t) V^\dagger W(t) V)$, so we have 
\begin{align}
    F(t)= 1 - \frac{1}{2d}Tr(C(t))=1-\frac{1}{2d}Tr\left[[W(t),V]^\dagger [W(t),V]\right]
\end{align}
\section{Equivalence of Dissipative Channels}
\label{equivalence_of_maps}
\subsection{Continuous Channel}
Notice that the instantaneous kicks in the evolution generated by (\ref{Flqop}) take place at integer steps and we can allow for continuous dephasing between the kicks. For $t \in (m,m+1)$ we have 
\begin{align}
    \dot{\rho} = -i[H_\text{I}, \rho] + \mathcal{D}(\rho) \equiv \mathcal{L}_{\text{cont}}(\rho).
\end{align}
Then at $t=m$, we apply the kick
\begin{align}
\rho(m^+)=\UK\rho(m^-)\UK^\dagger, \quad \UK=e^{-i \HK}
\end{align}
Therefore the dynamical map between an integer period is given by 
\begin{align}
    \rho_{m+1}=\UK \exp^{\mathcal{L}_\text{cont}}(\rho_m)\UK^\dagger .
\end{align}
where $\mathcal{L}_{\text{cont}}$ is the continuous channel Liouvillian. We can therefore define the continuous channel dynamical map as the composition of continuous time dissipation followed by a kick 
\begin{align}
    \mathcal{M}_{\text{cont}}(\rho):=\mathcal{U}_\text{K} \circ e^{\mathcal{L}_{\text{cont}}}(\rho).
\end{align}
\subsection{Stroboscopic Channel}
For the stroboscopic channel we simply allow for the dissipation to accumulate over the integer time-period and do not resolve the micro-time within that period and so the system density matrix evolves via the following map
\begin{align}
    \mathcal{U}_{\text{sys}}(\rho)&= U_{\text{sys}}[\bm{h}] \rho U^\dagger_{\text{sys}}[\bm{h}] \nonumber\\
\end{align}
We then introduce a dissipation channel $\mathcal{E}_{\text{diss}}$ that acts once per period, and the complete dynamical map looks like
\begin{align}
    \mathcal{M}_{\text{strob}}(\rho):=\mathcal{E
    }_{\text{diss}} \circ \mathcal{U}_{\text{sys}}(\rho).
\end{align}
Such a stroboscopic channel closely resembles experimental setups, where drives are practically instantaneous and periodic and one only cares about integer stroboscopic times where Markovianity ensures that the environment is featureless on the scale of a period. So this can be thought of as a coarse-grained Floquet quantum channel.
\newline
In the weak dissipation limit $(\gamma \ll 1)$ both the continuous and the stroboscopic channel are equivalent up to $\mathcal{O}(\gamma)$ corrections.
\\[4pt]
Let's look at the continuous map
\begin{align}
    \mathcal{M}_{\text{cont}}= \mathcal{U}_{\text{K}} \circ e^{\mathcal{L}_{\text{cont}}}
\end{align}
where
\begin{align}
    \mathcal{L}_\text{cont} &\sim -i[\HI,\rho]+\gamma \mathcal{D}(\rho) \nonumber
\end{align}
Using BCH
\begin{align}
    e^{-i[\HI,\cdot]+\gamma \mathcal{D}(\cdot)}&= e^{-i[\HI,\cdot]}e^{\gamma \mathcal{D}(\cdot)} \exp \left(-\frac{\gamma}{2}[-i[\HI,\cdot],\mathcal{D}]+\cdots\right) \nonumber \\
    &\sim e^{-i[\HI,\cdot]}e^{\gamma \mathcal{D}(\cdot)} + \mathcal{O}(\gamma)
\end{align}
Plugging this into the map we have
\begin{align}
    \mathcal{M}_{\text{cont}}=\mathcal{U}_{\text{K}} \circ e^{-i[\HI,\cdot]} \circ e^{\gamma \mathcal{D}(\cdot)} + \mathcal{O}(\gamma).
\end{align}
Since the composition ${U}_{\text{K}} \circ e^{-i[\HI,\cdot]}$ is exactly $\mathcal{U}_{\text{sys}}=\mathcal{U}_{\text{K}}\circ \mathcal{U}_{I}$, the continuous map is 
\begin{align}
    \mathcal{M}_{\text{cont}}=\mathcal{U}_{\text{sys}} \circ e^{\gamma \mathcal{D}(\cdot)} + \mathcal{O}(\gamma) \nonumber\\
    \mathcal{M}_{\text{cont}}=\mathcal{M}_{\text{strob}} + \mathcal{O}(\gamma).
\end{align}
This is the exact same form as $\mathcal{M}_{\text{strob}}$, up to where you put the dissipation channel before or after the unitary channel. To first order in $\gamma$, the choice only differs by $\mathcal{O}(\gamma)$.
\section{Logarithmic Negativity}
\label{app:Log-neg}
For a given bi-partition $A|A^c$ where $A=\{1,\cdots,N\}$ and $A^c=\{N+1,\cdots,L\}$, if we have a many-body global pure state $\rho_{AA^c}$ (flat spectrum), the Logarithmic Negativity (LN) is given by 
\begin{align}
    \mathcal{E}(\rho_{AA^c})\equiv\log||\rho_{AA^c}^{T}||_1,
\end{align}
where $T$ is the partial transpose on $A^c$ and $||\cdot||_1$ is the trace norm. Rényi entropy of order $\alpha$ for the reduced density matrix $\rho_A=\Tr_{A^c}(\rho_{AA^c})$ is given by 
\begin{align}
    S_A^{(\alpha)}= \frac{1}{1-\alpha}\log \Tr{(\rho_A^\alpha)}.
\end{align}
For $\alpha=1/2$, we have 
\begin{align}
    S_a^{1/2} = 2\log \Tr{(\rho_A^{1/2})}.
\end{align}
The Schmidt decomposition of the pure state is 
\begin{align}
    \ket{\psi}_{AA^c} = \sum_i\sqrt{\lambda_i}\ket{i}_A \otimes \ket{i}_{A^c},\quad\lambda_i \geq0,\quad\sum_i\lambda_i =1.
\end{align}
The reduce density matrix is given by 
\begin{align}
    \rho_A = \sum_i \lambda_i \ket{i}\bra{i}
\end{align}
so 
\begin{align}
    \Tr\rho_A^{1/2}=\sum_i \sqrt{\lambda_i}
\end{align}
and 
\begin{align}
    S^{1/2}_A = 2 \log \left(\sum_i \sqrt{\lambda_i}\right)
\end{align}
Now we compute $||\rho^T||_1$,
\begin{align}
    \rho_{AA^c}&= \sum_{ij}\sqrt{\lambda_i \lambda_j} \ket{i}\bra{j}_A \otimes \ket{i}\bra{j}_{A^c}\\
    \rho_{AA^c}^T&=\sum_{ij}\sqrt{\lambda_i \lambda_j} \ket{i}\bra{j}_A \otimes \ket{j}\bra{i}_{A^c}.
\end{align}
Calculating the trace norm
\begin{align}
    ||\rho_{AA^c}^T||_1&=\Tr\sqrt{\rho_{AA^c}^{T\dagger}\rho_{AA^c}^T}\\
    &= \sum_i \lambda_i + 2\ \sum_{i<j}\sqrt{\lambda_i \lambda_j}.
\end{align}
Note that $\sum_i \lambda_i =1$ so we have
\begin{align}
    \left(\sum_i \sqrt{\lambda_i}\right)^2 = \sum_i \lambda_i + 2\ \sum_{i<j}\sqrt{\lambda_i \lambda_j}.
\end{align}
Therefore, 
\begin{align}
    ||\rho_{AA^c}^T||_1&= \left(\sum_i \sqrt{\lambda_i}\right)^2
\end{align}
Taking the $\log$, 
\begin{align}
    \mathcal{E}(\rho_{AA^c})=2 \log \left(\sum_i \sqrt{\lambda_i}\right)=S^{1/2}_A
\end{align}
We also know that the spectrum of the reduced density matrix, in the thermodynamic limit, is flat, which means there are $2^{min(2t,N)}$ non-zero eigenvalues euqal to $2^{-\min(2t,N)}$ and rest are zero. So
\begin{align}
    \mathcal{E}=2\log\left(2^{\min(2t,N)/2}\right)=\min(2t,N)\log2.
\end{align}
% %\section{Equivalence of Dissipative Channels}
\label{equivalence_of_maps}
\subsection{Continuous Channel}
Notice that the instantaneous kicks in the evolution generated by (\ref{Flqop}) take place at integer steps and we can allow for continuous dephasing between the kicks. For $t \in (m,m+1)$ we have 
\begin{align}
    \dot{\rho} = -i[H_\text{I}, \rho] + \mathcal{D}(\rho) \equiv \mathcal{L}_{\text{cont}}(\rho).
\end{align}
Then at $t=m$, we apply the kick
\begin{align}
\rho(m^+)=\UK\rho(m^-)\UK^\dagger, \quad \UK=e^{-i \HK}
\end{align}
Therefore the dynamical map between an integer period is given by 
\begin{align}
    \rho_{m+1}=\UK \exp^{\mathcal{L}_\text{cont}}(\rho_m)\UK^\dagger .
\end{align}
where $\mathcal{L}_{\text{cont}}$ is the continuous channel Liouvillian. We can therefore define the continuous channel dynamical map as the composition of continuous time dissipation followed by a kick 
\begin{align}
    \mathcal{M}_{\text{cont}}(\rho):=\mathcal{U}_\text{K} \circ e^{\mathcal{L}_{\text{cont}}}(\rho).
\end{align}
\subsection{Stroboscopic Channel}
For the stroboscopic channel we simply allow for the dissipation to accumulate over the integer time-period and do not resolve the micro-time within that period and so the system density matrix evolves via the following map
\begin{align}
    \mathcal{U}_{\text{sys}}(\rho)&= U_{\text{sys}}[\bm{h}] \rho U^\dagger_{\text{sys}}[\bm{h}] \nonumber\\
\end{align}
We then introduce a dissipation channel $\mathcal{E}_{\text{diss}}$ that acts once per period, and the complete dynamical map looks like
\begin{align}
    \mathcal{M}_{\text{strob}}(\rho):=\mathcal{E
    }_{\text{diss}} \circ \mathcal{U}_{\text{sys}}(\rho).
\end{align}
Such a stroboscopic channel closely resembles experimental setups, where drives are practically instantaneous and periodic and one only cares about integer stroboscopic times where Markovianity ensures that the environment is featureless on the scale of a period. So this can be thought of as a coarse-grained Floquet quantum channel.
\newline
In the weak dissipation limit $(\gamma \ll 1)$ both the continuous and the stroboscopic channel are equivalent up to $\mathcal{O}(\gamma)$ corrections.
\\[4pt]
Let's look at the continuous map
\begin{align}
    \mathcal{M}_{\text{cont}}= \mathcal{U}_{\text{K}} \circ e^{\mathcal{L}_{\text{cont}}}
\end{align}
where
\begin{align}
    \mathcal{L}_\text{cont} &\sim -i[\HI,\rho]+\gamma \mathcal{D}(\rho) \nonumber
\end{align}
Using BCH
\begin{align}
    e^{-i[\HI,\cdot]+\gamma \mathcal{D}(\cdot)}&= e^{-i[\HI,\cdot]}e^{\gamma \mathcal{D}(\cdot)} \exp \left(-\frac{\gamma}{2}[-i[\HI,\cdot],\mathcal{D}]+\cdots\right) \nonumber \\
    &\sim e^{-i[\HI,\cdot]}e^{\gamma \mathcal{D}(\cdot)} + \mathcal{O}(\gamma)
\end{align}
Plugging this into the map we have
\begin{align}
    \mathcal{M}_{\text{cont}}=\mathcal{U}_{\text{K}} \circ e^{-i[\HI,\cdot]} \circ e^{\gamma \mathcal{D}(\cdot)} + \mathcal{O}(\gamma).
\end{align}
Since the composition ${U}_{\text{K}} \circ e^{-i[\HI,\cdot]}$ is exactly $\mathcal{U}_{\text{sys}}=\mathcal{U}_{\text{K}}\circ \mathcal{U}_{I}$, the continuous map is 
\begin{align}
    \mathcal{M}_{\text{cont}}=\mathcal{U}_{\text{sys}} \circ e^{\gamma \mathcal{D}(\cdot)} + \mathcal{O}(\gamma) \nonumber\\
    \mathcal{M}_{\text{cont}}=\mathcal{M}_{\text{strob}} + \mathcal{O}(\gamma).
\end{align}
This is the exact same form as $\mathcal{M}_{\text{strob}}$, up to where you put the dissipation channel before or after the unitary channel. To first order in $\gamma$, the choice only differs by $\mathcal{O}(\gamma)$.
% %\section{Logarithmic Negativity}
\label{app:Log-neg}
For a given bi-partition $A|A^c$ where $A=\{1,\cdots,N\}$ and $A^c=\{N+1,\cdots,L\}$, if we have a many-body global pure state $\rho_{AA^c}$ (flat spectrum), the Logarithmic Negativity (LN) is given by 
\begin{align}
    \mathcal{E}(\rho_{AA^c})\equiv\log||\rho_{AA^c}^{T}||_1,
\end{align}
where $T$ is the partial transpose on $A^c$ and $||\cdot||_1$ is the trace norm. Rényi entropy of order $\alpha$ for the reduced density matrix $\rho_A=\Tr_{A^c}(\rho_{AA^c})$ is given by 
\begin{align}
    S_A^{(\alpha)}= \frac{1}{1-\alpha}\log \Tr{(\rho_A^\alpha)}.
\end{align}
For $\alpha=1/2$, we have 
\begin{align}
    S_a^{1/2} = 2\log \Tr{(\rho_A^{1/2})}.
\end{align}
The Schmidt decomposition of the pure state is 
\begin{align}
    \ket{\psi}_{AA^c} = \sum_i\sqrt{\lambda_i}\ket{i}_A \otimes \ket{i}_{A^c},\quad\lambda_i \geq0,\quad\sum_i\lambda_i =1.
\end{align}
The reduce density matrix is given by 
\begin{align}
    \rho_A = \sum_i \lambda_i \ket{i}\bra{i}
\end{align}
so 
\begin{align}
    \Tr\rho_A^{1/2}=\sum_i \sqrt{\lambda_i}
\end{align}
and 
\begin{align}
    S^{1/2}_A = 2 \log \left(\sum_i \sqrt{\lambda_i}\right)
\end{align}
Now we compute $||\rho^T||_1$,
\begin{align}
    \rho_{AA^c}&= \sum_{ij}\sqrt{\lambda_i \lambda_j} \ket{i}\bra{j}_A \otimes \ket{i}\bra{j}_{A^c}\\
    \rho_{AA^c}^T&=\sum_{ij}\sqrt{\lambda_i \lambda_j} \ket{i}\bra{j}_A \otimes \ket{j}\bra{i}_{A^c}.
\end{align}
Calculating the trace norm
\begin{align}
    ||\rho_{AA^c}^T||_1&=\Tr\sqrt{\rho_{AA^c}^{T\dagger}\rho_{AA^c}^T}\\
    &= \sum_i \lambda_i + 2\ \sum_{i<j}\sqrt{\lambda_i \lambda_j}.
\end{align}
Note that $\sum_i \lambda_i =1$ so we have
\begin{align}
    \left(\sum_i \sqrt{\lambda_i}\right)^2 = \sum_i \lambda_i + 2\ \sum_{i<j}\sqrt{\lambda_i \lambda_j}.
\end{align}
Therefore, 
\begin{align}
    ||\rho_{AA^c}^T||_1&= \left(\sum_i \sqrt{\lambda_i}\right)^2
\end{align}
Taking the $\log$, 
\begin{align}
    \mathcal{E}(\rho_{AA^c})=2 \log \left(\sum_i \sqrt{\lambda_i}\right)=S^{1/2}_A
\end{align}
We also know that the spectrum of the reduced density matrix, in the thermodynamic limit, is flat, which means there are $2^{min(2t,N)}$ non-zero eigenvalues euqal to $2^{-\min(2t,N)}$ and rest are zero. So
\begin{align}
    \mathcal{E}=2\log\left(2^{\min(2t,N)/2}\right)=\min(2t,N)\log2.
\end{align}
% %\include{AppendixE}
% %\inclide{AppendixG}
\end{widetext}
\end{document}

